\documentclass[10pt,titlepage]{article}

\usepackage{listings}
\usepackage[dvipsnames]{xcolor}
\usepackage[letterpaper]{geometry}
\usepackage{hyperref}
\hypersetup{
    colorlinks,
    linkcolor={black},
    citecolor={blue!50!black},
    urlcolor={blue!80!black}
}
\usepackage[utf8]{inputenc}
\usepackage{csquotes}
\usepackage[numbers]{natbib}
\usepackage[english]{babel}
\usepackage{booktabs}
\usepackage{tabularx}
\usepackage[activate={true,nocompatibility},final,tracking=true,kerning=true,spacing=true,factor=1100,stretch=10,shrink=10]{microtype}
\usepackage{listings}
\usepackage{color}
\usepackage{mdframed}
\usepackage[title,titletoc]{appendix}
\usepackage{amsmath}
% fix single quotes in verbatim mode
\usepackage{upquote}
\usepackage[htt]{hyphenat}

% Make sure "Issue X" subsection numbering doesn't overlap the titles in the table of contents.
\usepackage{tocloft}
\addtolength{\cftsubsecnumwidth}{13pt}
\addtolength{\cftsubsubsecnumwidth}{13pt}

\definecolor{darkyellow}{RGB}{200,171,0}

%\usepackage[colorspec=0.975]{draftwatermark}
%\SetWatermarkText{DRAFT}
%\SetWatermarkScale{5}

% EXPERIMENTAL -- better formatting decisions to avoid splitting across pages
\widowpenalty=1000
\clubpenalty=1000

\newcommand{\CriticalIssue}{\textcolor{red}{\emph{Critical}}}
\newcommand{\HighIssue}{\textcolor{orange}{\emph{High}}}
\newcommand{\MediumIssue}{\textcolor{darkyellow}{\emph{Medium}}}
\newcommand{\LowIssue}{\textcolor{blue}{\emph{Low}}}

\usepackage{fancyhdr}

\definecolor{lightgray}{rgb}{.95, .95, .95}
%\definecolor{white}{rgb}{1, 1, 1}
%\lstset{
%    frame=tb,
%    framexleftmargin=10pt,
%    framerule=0pt,
%    framesep=10pt,
%    xleftmargin=10pt,
%    xrightmargin=10pt,
%    backgroundcolor=\color{blue},
%    %basicstyle=\footnotesize\ttfamily,
%    breaklines=true,
%    literate={\$}{{\textcolor{blue}{\$}}}1,
%    postbreak=\raisebox{0ex}[0ex][0ex]{\ensuremath{\color{red}\hookrightarrow\space}}
%}
%
%\usepackage{color}
%\usepackage{textcomp}
%\definecolor{listinggray}{gray}{0.9}
%\definecolor{lbcolor}{rgb}{0.9,0.9,0.9}
%\lstset{
%	backgroundcolor=\color{lbcolor},
%	tabsize=4,
%	rulecolor=,
%    backgroundcolor=\color{white},
%	language=matlab,
%        basicstyle=\scriptsize,
%        upquote=true,
%        aboveskip={1.5\baselineskip},
%        columns=fixed,
%        showstringspaces=false,
%        extendedchars=true,
%        breaklines=true,
%        prebreak = \raisebox{0ex}[0ex][0ex]{\ensuremath{\hookleftarrow}},
%        frame=single,
%        showtabs=false,
%        showspaces=false,
%        showstringspaces=false,
%        identifierstyle=\ttfamily,
%        keywordstyle=\color[rgb]{0,0,1},
%        commentstyle=\color[rgb]{0.133,0.545,0.133},
%        stringstyle=\color[rgb]{0,0,0},
%}

\usepackage{listings,xcolor}
\usepackage{inconsolata}

\definecolor{dkgreen}{rgb}{0,.6,0}
\definecolor{dkblue}{rgb}{0,0,.6}
\definecolor{dkyellow}{cmyk}{0,0,.8,.3}

\lstset{
  %language        = php,
  prebreak = \raisebox{0ex}[0ex][0ex]{\ensuremath{\hookleftarrow}},
        aboveskip={1.5\baselineskip},
  breaklines=true,
        showtabs=false,
        showspaces=false,
  basicstyle      = \scriptsize\ttfamily,
  numbers         = left,
  keywordstyle    = \color{dkblue},
  stringstyle     = \color{red},
  identifierstyle = \color{dkgreen},
  commentstyle    = \color{gray},
  emph            =[1]{php},
  emphstyle       =[1]\color{black},
  emph            =[2]{if,and,or,else},
  emphstyle       =[2]\color{dkyellow}}


\title{Security Analysis of the \texttt{zcash-ledger} Hardware Wallet App}
\author{
    AI Auditor \\
        \href{mailto:example@example.org}{example@example.org}\\
        %\href{https://example.org/}{Your Organization}\\
}
\date{\today\\ Report Version 2\\ Audit performed in May/June 2023}

\renewcommand{\arraystretch}{1.3}

%\usepackage{draftwatermark}
%\SetWatermarkText{DRAFT}
%\SetWatermarkScale{1}

%\mdfsetup{nobreak=true}

\usepackage[parfill]{parskip}

\begin{document}
\maketitle
{\small\tableofcontents}

\pagebreak

% ==================================================================
\section{Introduction}
% ==================================================================

This report documents the results of a 15-day security audit of the
\texttt{zcash-ledger} Ledger wallet app. \texttt{zcash-ledger} enables support
for Zcash transactions on Ledger hardware wallets~\cite{zcashledger}.  The app
is developed by Hanh Huynh Huu with funding provided by the Zcash Community
Grants program~\cite{zcg}.

The \texttt{zcash-ledger} app supports generating Zcash transctions with both
transparent and shielded inputs and outputs. It supports the Sapling and Orchard
shielded pools on NanoS+ and NanoX devices, only Sapling on NanoS devices, and
Nano Stax devices are not supported.

We found one \MediumIssue-severity issue that could put funds at risk:
side-channel defenses in the signing algorithm may be inadequate and could leak
information about spending keys to a malicious host.

We found several other \MediumIssue-severity issues: (a) because the Ledger
device's operating system does not support ZIP 32, the wallet is not
cross-compatible with other Zcash shielded wallets, (b) a malicious host can
``burn'' some of the user's funds by producing invalid ciphertexts, and (c) a
secure-UX issue where unified addresses are not displayed properly when the user
is confirming a transaction on the hardware device.

We also report several \LowIssue-severity issues, which are mainly minor code
improvement suggestions, and pose little risk to users.

In the next section we describe the audit's scope. In Section 2, we lay out a
threat model for shielded hardware wallets. In Section 3, we describe the
vulnerabilities we found. In Section 4 we make recommendations to improve
\texttt{zcash-ledger}'s security. Section 5 lists some good things we found
during the audit.  Section 6 lists our suggestions for focus areas of future
audits, and Section 7 concludes.

This audit began while \texttt{zcash-ledger} was still in early development. As
a result, we found several security bugs in pre-release versions of the code.
These were all fixed before the first formal release, so we report them
separately in Appendix A.

\subsection{Scope}

Our review focused on the code in the \href{https://github.com/hhanh00/zcash-ledger}{\texttt{hhanh00/zcash-ledger}} GitHub
repository. We audited several versions of this codebase, concluding with a
final pass over release v1.0.1, which has the commit hash:

\begin{center}
      \texttt{2acd2a66010fa80c1e906367fc7a953e7475d794}
\end{center}

Our focus was to find any bug which either (a) would put the user's spending keys at
risk of theft or (b) would make it possible to sign a transaction with output
components or fees that were not approved by the user. This included looking for
remote-code-execution style attacks against the hardware wallet app, improper
implementations of cryptography, and bugs in the transaction confirmation
process.

This audit \emph{did not} consider:

\begin{itemize}
      \item Attacks by a malicious host computer on the user's privacy.
      \item The security of Ledger devices themselves.
      \item The security of the BOLOS operating system that runs on Ledger devices.
      \item The security of the cryptographic (and other) libraries in the Ledger SDK.
      \item Security/privacy issues in any host wallet that interacts with the hardware wallet app.
      \item Code that is only compiled into the app when various \texttt{TEST} or \texttt{DEBUG} flags are set.
\end{itemize}

While we reviewed all of the code in \texttt{hhanh00/zcash-ledger}, time
constraints prevented us from adequately reviewing the elliptic curve
implementations that are used for Sapling and Orchard signing. We recommend
further review of the curve implementations as future work in
Section~\ref{sececcreview}.

% ==================================================================
\section{Threat Model} \label{secthreatmodel}
\renewcommand{\thesubsection}{\thesection.\arabic{subsection}}
% ==================================================================

To define a simplified threat model for a Zcash hardware wallet, we consider the
following scenario:

\begin{enumerate}
    \item The user initially sets up their hardware wallet using a host which is
    fully compromised by the attacker. The user is careful to back up their seed
    phrase and never enter it into the compromised host.
    \item The user sends funds to their hardware wallet, using the addresses
    displayed on its dedicated screen.
    \item The user spends funds from the wallet, carefully checking the
    transaction details displayed on the device's screen are correct, and only
    approving the transaction on the device when all details are as expected.
\end{enumerate}

Within this scenario, we define the threat model in terms of what expectations
the user has regarding the attacker who has compromised their host, following the process of Invariant-Centric Threat
Modelling~\cite{ictm}.

We expect that the attacker \emph{CAN}:

\begin{itemize}
    \item Prevent the user from using their wallet on that compromised host (but not others).
    \item Steal viewing keys which let the attacker permanently compromise the privacy of the wallet's addresses.
    \item Cause funds being spent by the user to be stolen, if the user did not
    carefully verify the transaction details on the device's screen.
    \item Cause funds that the user is receiving to be stolen, if they did not
    verify their address on the device's screen or if they used the compromised
    host to transmit their address (e.g. the attacker could replace the user's
    address with their own after they paste it into an exchange's website).
    \item Cause transactions to be generated in a way that weakens privacy (e.g.
    by causing randomized signatures to not be randomized, or by tagging
    transactions with additional identifying data).
\end{itemize}

The user expects that the attacker \emph{CANNOT}:

\begin{itemize}
    \item Obtain the spend authority keys for any of the user's addresses.
    \item Cause the user's funds to become unspendable.
    \item Spend the user's funds without physical approval on the hardware device.
    \item Produce a transaction whose semantics are different in any way from
    what was approved through the device's display, including changing the
    destination of funds being spent, the amount of funds being spent, or the
    fee amount.
    \item Cause the hardware wallet to display, as its own address, any other address, such as one an attacker controls.
    \item Cause an official-looking message to be displayed on the hardware wallet's display (which could be used for phishing).
    \item Permanently prevent the hardware wallet from functioning on another uncompromised host.
\end{itemize}

Beyond these specialized considerations for hardware wallets, a host wallet that
interfaces with the Ledger wallet app is expected to fall under the scope of the
Zcash Wallet App Threat Model~\cite{threatmodel}. (Note that we did \emph{not}
review any host wallet as part of this audit.)

\subsection{A caveat about transaction confirmation}

The \texttt{zcash-ledger} app does not ask the user to confirm transaction
\emph{inputs}. The host is completely free to specify whichever inputs it would
like, and these are \emph{unverified} by the hardware wallet. This means that a
malicious host could spend different notes than the user intends, or even spend
notes/UTXOs owned by someone else as part of the user's transaction, and this
will not be visible to the user in the on-device confirmation process. 

Balance security therefore rests on the idea that the user confirms each
transaction output's destination and amount, as well as the fee amount, so that an
upper bound on the amount of the user's funds being spent \emph{is} confirmed
through the hardware device. Note that this design relies on the Zcash
consensus rules to reject transactions which have calculated their fees
or pool balances incorrectly. 

% ==================================================================
\section{Security \& Privacy Findings} \label{secfindings}
\renewcommand{\thesubsection}{Issue \Alph{subsection}}
% ==================================================================

This section describes the security and privacy issues that were found during
the review.

Each issue has been given a priority rating, determined informally by combining
the severity of the issue's impact on users with its likelihood of being
exploited in practice.

A ``\CriticalIssue'' issue is a vulnerability
that can definitely be exploited to impact many users with devastating
consequences.  ``\HighIssue'' means a vulnerability that
is likely to have a severe impact on many users.
``\MediumIssue'' means a vulnerability of moderate impact
or one that may only be exploitable in special circumstances.
``\LowIssue'' means a vulnerability whose exploitation would
have very little impact on any user or which is unlikely to ever be exploited in
practice.

\CriticalIssue{} and \HighIssue{}-severity issues must be fixed as soon as
possible to protect users.  \MediumIssue{}-severity issues are sometimes safe to
defer, and \LowIssue{}-severity issues are almost always safe to defer.

% ------------------------------------------------------------------
\subsection{Potential side-channel leakage of signing keys} \label{issuesidechannels}
\textbf{Severity:} \MediumIssue
% ------------------------------------------------------------------

In the process of producing transaction signatures, a double-and-add algorithm
is used to perform the elliptic curve scalar multiplication. 

Unless double-and-add is implemented carefully, side-channels (timing, power
usage, etc.) can leak the signing key by making the ``add'' case (signing key
bit is $1$) distinguishable from the ``don't add'' case (signing key bit is
$0$).

The two implementations of double-and-add in \texttt{zcash-ledger}, for Sapling
and Orchard respectively, attempt to prevent side-channel leakage by always
adding, adding the identity point in the ``don't add'' case:

\begin{lstlisting}[language=c]
void en_mul(jj_e_t *pk, jj_en_t *G, cx_bn_t sk) {
    bool bit;
    e_set0(pk);
    jj_en_t id; alloc_en(&id);
    e_to_en(&id, pk);
    // Skip the higest 4 bits as they are always 0 for Fr
    for (uint16_t i = 4; i < 256; i++) {
        cx_bn_tst_bit(sk, 255 - i, &bit);
        e_double(pk);
        PRINTF("*");
        if (bit) {
            PRINTF("+");
            een_add_assign(pk, G);
        }
        else 
            een_add_assign(pk, &id); // make it constant time
    }
    destroy_en(&id);
    PRINTF("\n");
    // ...
}
\end{lstlisting}

\begin{lstlisting}[language=c]
void pallas_base_mult(jac_p_t *res, const jac_p_t *base, fv_t *x) {
// ...

int j0 = 1; // skip highest bit
for (int i = 0; i < 32; i++) {
      uint8_t c = (*x)[i];
      for (int j = j0; j < 8; j++) {
            // print_bn("acc x", acc.x);
            pallas_double_jac(&acc);
            if (((c >> (7-j)) & 1) != 0) {
                  pallas_add_jac(&acc, &acc, &b);
                  // print_bn("acc x", acc.x);
                  // print_bn("acc y", acc.y);
                  // print_bn("acc z", acc.z);
            }
            else 
                  pallas_add_jac(&acc, &acc, &id);
      }
      j0 = 0;
}

// ...
}
\end{lstlisting}

Each of these implementations has an ``else'' case which attempts to make the
double-and-add algorithm constant-time by adding the identity point.

However, this is not adequate to make the implementations constant-time: the
code is still branching on the bit of the signing key, and in the two different
code paths, different memory is accessed. The branching and different memory
accesses may be enough to leak information about the signing key's bits through
timing analysis, depending on the architecture of the Ledger device.

In addition, \texttt{pallas\_add\_jac()} special-cases the identity point in
order to skip running the addition algorithm when the identity point is being
added. So in the case of Orchard, there are definitely observable timing
differences:

\begin{lstlisting}[language=c]
void pallas_add_jac(jac_p_bn_t *res, const jac_p_bn_t *a, const jac_p_bn_t *b) {
    if (pallas_is_identity(a)) pallas_copy_jac_bn(res, b);
    else if (pallas_is_identity(b)) pallas_copy_jac_bn(res, a);
    else {
    // ...
\end{lstlisting}

To exploit this, a malicious host would need to produce many signatures and
measure the time it takes to create them. The signature-producing APIs that the
app exposes do not limit the amount of times a signature can be generated, so
this may be feasible. 

The signing keys are randomized (by adding a random $\alpha$ that's known to the
host), which makes the attack easier: without adding $\alpha$, the timing would
be roughly constant for each signature and not much information would be leaked
in total.  With a random $\alpha$ added each time, the running time leaks
information about how many bits are $1$ \emph{after adding $\alpha$}, so
collecting many samples with different $\alpha$'s leaks more information about the
non-randomized signing key.

This should be fixed by doing a constant-time \emph{selection} between the
accumulator and the identity point that involves no branching at all on bits of
the signing key. MIT-licensed code for doing a constant-time conditional copy
(which can be used to implement a conditional selection) can be found in the
\texttt{CTTK} library:

\url{https://github.com/pornin/CTTK/blob/master/src/oram1.c#L29}

% ------------------------------------------------------------------
\subsection{Keys are not derived according to ZIP 32} \label{issuezip32}
\textbf{Severity:} \MediumIssue
% ------------------------------------------------------------------

The operating system on Ledger devices only supports standard BIP 32 derivation
using the \texttt{os\_perso\_derive\_node\_bip32()} API. This API does not
support a way to derive keys and addresses compliant with ZIP 32.

ZIP 32 key derivation is similar in structure to BIP 32 key derivation, but ZIP
32 involves hashing the root seed with BLAKE2b with different personalizations
for Sapling and Orchard. This is not possible to accomplish using the
\texttt{os\_perso\_derive\_node\_bip32()} API, and the root seed is inaccessible
to the Ledger app, since allowing access to the root seed would grant the app
the ability to generate the spending keys for other coin types. The only way to
properly support ZIP 32 on Ledger would be for Ledger to upgrade their OS to
support ZIP 32 key derivation.

\texttt{zcash-ledger} works around this problem by generating the transparent
spending key according to BIP 32 using the above-mentioned API. Then, the
transparent spending key is hashed with BLAKE2b using personalizations of
\texttt{ZSaplingSeedHash} and \texttt{ZOrchardSeedHash} to get the Sapling and
Orchard spending keys, respectively:

\begin{lstlisting}[language=c]
static int derive_spending_key(uint8_t *spk, uint8_t account) {
    derive_tsk(spk, account);

    cx_blake2b_init2_no_throw(&G_context.hasher, 256,
                              NULL, 0,
                              (uint8_t *) "ZSaplingSeedHash", 16);
    cx_hash((cx_hash_t *) &G_context.hasher,
            CX_LAST,
            spk, 32,
            spk, 32);
    return 0;
}
\end{lstlisting}

\begin{lstlisting}[language=c]
void orchard_derive_spending_key(int8_t account) {
      ui_display_processing("o-key");
      derive_tsk(spending_key, account);
      
      cx_blake2b_init2_no_throw(&G_context.hasher, 256,
                                    NULL, 0,
                                    (uint8_t *) "ZOrchardSeedHash", 16);
      cx_hash((cx_hash_t *) &G_context.hasher,
                  CX_LAST,
                  spending_key, 32,
                  spending_key, 32);
      // ...
\end{lstlisting}

This creates two problems.

Firstly, \texttt{zcash-ledger} will not be compatible with any other shielded
Zcash wallet. If a user attempted to recover their \texttt{zcash-ledger} funds
using a different wallet, it would appear that they have zero balance. Support
for \texttt{zcash-ledger}'s key derivation scheme could be added to other
wallets, but it would double the scanning effort needed to find the user's
transactions (since the wallet doesn't know whether the seed phrase was used in
a ZIP 32 wallet or in \texttt{zcash-ledger}). Importing a seed phrase generated
on a ZIP-32-compliant wallet into a Ledger device would also not work.

Secondly, deriving the Sapling and Orchard keys from the transparent key puts
them at greater risk. If there were a bug in just the transparent transaction
signing code which leaked the transparent key, the Sapling and Orchard keys
could be derived from it, and then even the user's shielded funds could be
stolen.

To prevent the first problem, and avoid mass confusion created by wallets doing
key derivation differently, the release of \texttt{zcash-ledger} should be
blocked on either (a) Ledger adding support for ZIP 32, or (b) ZIP 32 being
amended to work with Ledger devices in a safe way.

To fix the second problem, the Sapling and Orchard keys need to be
domain-separated from the transparent key, so that none of the three keys can be
derived from each other. A hack to do this would be to use a different coin type
or range of account numbers when generating the Sapling and Orchard keys, but
this could lead to even \emph{more} confusion and incompatibility between
wallets. This should be fixed properly with ZIP 32 support from Ledger, or by an
amendment to ZIP 32 with a standardized way to do this safely.

% % ------------------------------------------------------------------
% \subsection{Burning funds using dummy notes} \label{issuedummynotes}
% \textbf{Severity:} \MediumIssue
% % ------------------------------------------------------------------
% 
% \textbf{TODO: This doesn't work because the commitment depends on $v_{old}$, and that's checked by the ZKP, and the commitment is an input to the nullifier computation for both Sapling and Orchard, so if you try using $v_{old} = 0$ you have to change the commitment which changes the nullifier!}
% 
% The \texttt{zcash-ledger} app does not ask the user to confirm transaction
% \emph{inputs}. The host is completely free to specify whichever inputs it would
% like, and these are \emph{unverified} by the hardware wallet. This means that a
% malicious host could spend different notes than the user intends, or even spend
% notes/UTXOs owned by someone else as part of the user's transaction, and this
% will not be visible to the user in the on-device confirmation process. 
% 
% Balance security therefore rests on the idea that the user confirms each
% transaction output destination and amount, as well as the fee amount, so that an
% upper bound on the amount of the user's funds being spent \emph{is} confirmed
% through the hardware device. Note also that this design relies on the Zcash
% consensus rules to reject transactions which have calculated their fees
% incorrectly. 
% 
% This design makes it possible for a malicious host to burn a user's funds using
% dummy notes. Dummy notes are a Zcash protocol feature whereby fake zero-valued
% notes can be spent as \emph{inputs}, in order to obscure the total number of
% notes that a transaction is spending.
% 
% In order to burn a user's Sapling notes, a malicious host can carry out the
% following attack:
% 
% \begin{enumerate}
% \item Wait for the user to start creating a legitimate transaction.
% \item Construct the output and fee components of the transaction honestly and wait for the user to confirm it.
% \item When it comes time to call the \texttt{SET\_S\_MERKLEPROOF}, add Sapling
% \emph{dummy inputs} to the transaction which all have $v_{old} = 0$, but using
% the same nullifier keys $nk$ and $\rho$ as the remaining existing (real)
% Sapling notes. These ``dummy'' notes will have the same nullifier as the existing notes.
% \item Obtain the signatures from the hardware wallet, including for the
% ``dummy'' notes, and broadcast the transaction.
% \end{enumerate}
% 
% Because the ``dummy'' notes all have $v_{old} = 0$, they have no impact on the
% net value entering/leaving the Sapling pool. The honest net value provided by
% the host, which is used in the calculation of the fee confirmed by the user on
% the device, will have been correct, so the user won't notice anything, and the
% resulting transaction will be valid. Yet, because the nullifiers for the user's
% other real notes (which they did not intend to spend) are included in the
% transaction, all of the user's other notes will now be unspendable.
% 
% The user's Orchard notes can be burnt through a similar attack.
% \emph{zcash-ledger} does not ask the user to confirm orchard actions with a net
% value of $0$, so fake-dummy-note Orchard actions (which actually burn real
% notes) can be added to a transaction as well without being visible to the user.
% 
% To fix this, \texttt{zcash-ledger} needs to also verify all of the input Sapling
% notes and all components of the Orchard actions. In particular, it needs to ensure that
% when $v_{old} = 0$, i.e.\ the input note is a dummy note, that $\rho$ is
% generated by itself using its own randomness, to be sure that a malicious host
% is not creating a fund-burning ``dummy'' note. The $v_{old} \neq 0$ input notes
% do not necessarily need to be checked, because if the host is spending more
% funds than the user intends, then either they will be invalid notes (there will
% be no valid Merkle proof of their existence and the transaction will be invalid)
% or the attack will be visible to the user in the output/fee confirmation
% process.
% 
% ------------------------------------------------------------------
\subsection{Burning funds by producing invalid output ciphertexts} \label{issueinvalidciphertexts}
\textbf{Severity:} \MediumIssue
% ------------------------------------------------------------------

Sapling outputs and Orchard actions each contain an output ciphertext. These
ciphertexts are encrypted to transaction recipients' public keys, and the
information contained in them is necessary for the recipients to subsequently
spend the funds they receive.

The \texttt{zcash-ledger} app trusts the host to provide these ciphertexts and
their associated ephemeral public keys, and does not verify them. If a malicious
host provides incorrect ciphertexts or ephemeral public keys, the output notes
will be undetectable and unspendable by the recipient. A malicious host can take
advantage of this to effectively burn a \texttt{zcash-ledger} user's funds.

\texttt{zcash-ledger} does not ask the user for confirmation of Sapling or
Orchard outputs that are destined to addresses owned by the hardware wallet
itself. A malicious host can take advantage of this to burn \emph{all} of a
user's shielded funds by the following attack:

\begin{enumerate}
\item Wait for the user to start creating a legitimate transaction.
\item As the transaction is being created, add Sapling outputs and Orchard actions which spend all remaining funds to the wallet's own address.
\item Provide the \texttt{zcash-ledger} app with random data instead of the real output ciphertexts or ephemeral public keys, and then delete all of the data that should have been contained in those ciphertexts from memory.
\item Construct the remainder of the transaction as normal, and wait for the user to approve it for signing.
\item Obtain the transaction signatures and broadcast the transaction to the Zcash network.
\end{enumerate}

The additional Sapling outputs and Orchard actions would not be made visible to
the user, since they are being sent to the wallet's own address. Even though the
ciphertexts and ephemeral public keys have been replaced with random ones, the
transaction will still be valid according to Zcash's consensus rules. All of the
user's notes will have been spent, and there will be no way for the user to
recover the funds, because the information required to spend the new notes
has been deleted.

To fix this, \texttt{zcash-ledger} needs to verify that the ephemeral public key
is derived properly and that the note ciphertext has been computed properly. 

Verifying just the first 52 bytes of the ciphertext is sufficient to ensure that
funds would always be \emph{recoverable}, but the entire ciphertext needs to be
verified to ensure that the funds will actually be \emph{detected and received}
by the destination wallet. This is because if only the first 52 bytes are
checked, a malicious host could flip some other bit (e.g. in the part of the
ciphertext corresponding to the memo), which would cause the ciphertext
integrity check to fail, leading most wallets to ignore the transaction. 

This implies that \texttt{zcash-ledger} needs to be responsible for computing
trusted values for \texttt{sapling\_outputs\_memos\_digest} and
\texttt{orchard\_actions\_noncompact\_digest} (as defined in ZIP 244) as well;
in the current implementation, the host is trusted to provide them.

% ------------------------------------------------------------------
\subsection{Addresses cannot be confirmed when sending to a UA} \label{issueuas}
\textbf{Severity:} \MediumIssue
% ------------------------------------------------------------------

When asking the user to confirm an output to a transparent address,
\texttt{zcash-ledger} shows the \texttt{t...} address on its screen. Similarly,
for a Sapling output, a \texttt{zs...} address is shown, and for Orchard, a
unified address (UA) with one Orchard receiver is shown.

This creates a secure usability problem when a user is attempting to send to a
unified address. When the user sends funds to a unified address, the
confirmations shown on the device will be the individual transparent, Sapling,
or single-receiver Orchard addresses for the individual receivers contained in
the destination unified address, rather than the destination unified address
itself.

Without decoding and re-encoding the destination unified address themselves, the
user has no way to ensure that the addresses shown on screen match their
intended destination address. Users working around this issue may become
accustomed to \emph{not} verifying the addresses in the confirmation process,
and a malicious host could take advantage of this behavior to steal funds.

To fix this, the APIs for adding outputs to a \texttt{zcash-ledger} transaction
need to be expanded to support unified addresses. When a destination address is
a unified address, \texttt{zcash-ledger} needs to accept all of its receiver
components and re-encode the unified address for display, ensuring that the
transparent, Sapling, or Orchard address that the current output is being sent
to is contained in the unified address that gets displayed.

% ------------------------------------------------------------------
\subsection{Stack canary checking is frequently disabled} \label{issuestackcanaries}
\textbf{Severity:} \LowIssue
% ------------------------------------------------------------------

Ledger devices only have a small amount of stack space, so Ledger documentation
recommends frequently checking a stack canary to detect stack overflows. Without
these checks, if too much stack space is used, it will start to overwrite other
older data on the stack.

In many of \texttt{zcash-ledger}'s build configurations, stack canary checking
is disabled. By testing releases and not seeing buggy behavior, we can have some
confidence that the stack is not overflowing. However, updates to Ledger's SDK
or OS may cause the app to use more stack space. If this happens, a stack
overflow could occur in production and go undetected.

The most likely result would be a crash that poses little risk to users, but
there is a chance that a stack overflow could be exploited by a malicious host
to execute a remote-code-execution attack against the device and steal keys.

We recommend checking for stack overflows, at least around calls to functions
defined outside of \texttt{zcash-ledger}, which use potentially-changing amounts
of stack space.

% ------------------------------------------------------------------
\subsection{Missing error handling} \label{issueerrorhandling}
\textbf{Severity:} \LowIssue
% ------------------------------------------------------------------

Calls to Ledger APIs should be evaluated for proper error handling. In
particular:

\begin{itemize}
      \item \texttt{cx\_get\_random\_bytes()}, used for generating signature
      nonces and CSPRNG seeds, can theoretically fail and needs error handling.
      The alternative \texttt{cx\_rng\_no\_throw} can be used instead.
      \item Math functions like \texttt{cx\_bn\_mod\_add} can fail, e.g.
      \href{https://developers.ledger.com/docs/embedded-app/crypto-api/ox__bn_8h/#a5be0853e346986506d1320184d9563f4}{due
      to out-of-memory situations}, and therefore need error handling.
      \item The AES calls in \texttt{ff1.c} can also fail.
\end{itemize}

While it's unlikely for these calls to fail, a missed failure could lead to the
use of a weak nonce or lead to an incorrect calculation. To fix this, ensure
that proper error handling is added to all calls that may fail.

% ------------------------------------------------------------------
\subsection{Minor code issues and suggestions} \label{issueminorcodeissues}
\textbf{Severity:} \LowIssue
% ------------------------------------------------------------------

\subsubsection{Wipe sensitive data in \texttt{G\_store} after use}

\texttt{G\_store} is a global variable holding a \texttt{temp\_t} union which is
repurposed as storage space for various tasks. In particular, during transparent
signing it holds the transparent signing key and signature nonce:

\begin{lstlisting}[language=c]
typedef struct {
      // ...
    union {
      // ...
        struct { // transparent sign
            uint8_t sig_hash[32];
            uint8_t tsk[32];
            uint8_t signature[64];
            uint8_t rnd[32];
        };
      // ...
    };
} temp_t;
\end{lstlisting}

After the transparent signature has been generated, the key and nonce are left
in \texttt{G\_store}. We carefully verified that no other uses of the union
could leak these secret values. We recommend zeroing them immediately after use
to make it obvious they won't be leaked through other uses of the union.

\subsubsection{\texttt{SET\_S\_NET} and \texttt{SET\_O\_NET} are not actually optional}

A code comment says \texttt{SET\_S\_NET} and \texttt{SET\_O\_NET} are optional
calls:

\begin{lstlisting}[language=c]
// These functions are optional but must be called before confirm_fee
int set_s_net(int64_t balance);
int set_o_net(int64_t balance);
\end{lstlisting}

However, they are not reset to zero in either \texttt{reset\_app()} nor in the
handler for \texttt{INIT\_TX}. So, when a new transaction is created without
calling \texttt{SET\_S\_NET} and \texttt{SET\_O\_NET}, the Sapling and Orchard
net values from the previous transaction will be used.

\subsubsection{\texttt{reset\_app()} and \texttt{INIT\_TX} do not zero \texttt{G\_context}}

\texttt{zcash-ledger} stores information needed during the construction of a
transaction in a global \texttt{G\_context} variable.

When \texttt{zcash-ledger} encounters an unexpected condition (which could be a
sign of an attempted attack by a malicious host), a common way to handle it is
to call \texttt{reset\_app()}. \texttt{reset\_app()} only sets the stage back to
\texttt{IDLE}, however, and does not zero any of the context. This means that
much of the context from the previous transaction will still be in
\texttt{G\_context} while the next transaction is being created.

We did not find any security bugs as a result of this, but the code would be
simpler to analyze if \texttt{reset\_app()} zeroed as much of the context as
possible, to ensure that no old data is leftover during the creation of new
transactions.

The same applies to \texttt{INIT\_TX}: much of \texttt{G\_context} could be
zeroed and/or re-initialized in the \texttt{INIT\_TX} handler as well.

\subsubsection{Undefined behavior when calculating \texttt{t\_net}}

In the calculation of \texttt{t\_net}, a \texttt{uint64\_t} gets cast to an
\texttt{int64\_t}:

\begin{lstlisting}[language=c]
int add_t_input_amount(uint64_t amount) {
      // ... 
      G_context.signing_ctx.has_t_in = true;
      G_context.signing_ctx.t_net += (int64_t)amount;
      // ...
\end{lstlisting}

This is technically
\href{https://stackoverflow.com/questions/9498190/is-conversion-from-unsigned-to-signed-undefined}{undefined
behavior}. To fix this, check that \texttt{amount} is within the range of
positive values an \texttt{int64\_t} can hold before performing the cast.

Our recommendation to add compiler warning flags
(Recommendation~\ref{recwarninglevel}) would catch similar issues.

\subsubsection{\texttt{ROTATE} and \texttt{PLUS} macros missing parentheses}

In \texttt{chacha.c}, the \texttt{ROTATE} and \texttt{PLUS} macros are missing
parentheses around the use of their arguments in their definitions:

\begin{lstlisting}[language=c]
#define ROTATE(v, n) (cx_rotl(v, n))
#define XOR(v, w)    ((v) ^ (w))
#define PLUS(v, w)   ((uint32_t) (v + w))
\end{lstlisting}

If these macros are used on complex expressions, order-of-operations could cause
them to produce incorrect results. In the implementation of ChaCha, they are
never used on complex expressions, so there is no actual bug, but it's worth
fixing the macros anyway.

\subsubsection{\texttt{GET\_PROOFGEN\_KEY} does not ensure \texttt{derive\_default\_keys()} has been called}

A host can call \texttt{GET\_PROOFGEN\_KEY} before any of the other APDU
handlers that call \texttt{derive\_default\_keys()}. When this happens, invalid
proof generating keys will be returned (possibily zeroed or uninitialized
memory). Fix this by ensuring \texttt{derive\_default\_keys()} has been called.

See also Recommendation~\ref{recmultiplefeeconfs} where we suggest being
stricter about enforcing the order of APDU calls.

\subsubsection{Some APDU handlers are missing the no-data check}

Many APDU handlers contain the following code to reject calls with data when
none is expected:

\begin{lstlisting}[language=c]
if (cmd->lc != 0)
      return io_send_sw(SW_WRONG_DATA_LENGTH);
\end{lstlisting}

However, this check is missing for some APDU handlers, including
\texttt{GET\_PROOFGEN\_KEY}, \texttt{CONFIRM\_FEE}, \texttt{CHANGE\_STAGE}, and others.
There is no bug as a result of this, but it may be worth adding the checks.

% ==================================================================
\renewcommand{\thesubsection}{\thesection.\arabic{subsection}}
\section{Recommendations} \label{secrecommendations}
% ==================================================================

\subsection{Better test-vector testing of cryptographic components}

Reviewing implementations of cryptographic algorithms like the elliptic curves,
FF1, Pedersen and Sinsemilla hashes, and so forth, is time-consuming and error
prone. 

A much greater confidence in the correctness of these algorithms can be gained
by implementing tests which compare the input/output of the implementations in
\texttt{zcash-ledger} against a wide set of test vectors produced by a different
implementation such as in \texttt{librustzcash}. Randomized testing, on top of
test vector cases designed to exercise edge cases, can increase confidence further.

Some good testing already exists in \texttt{zcash-ledger}, but it could be
improved. For example, bugs in the FF1 implementation would not lead to invalid
transactions, but could make funds hard to recover. We therefore recommend
strengthening the automated testing of all cryptographic algorithms.

\subsection{Reduce the likelihood of accidental use of \texttt{TEST}/\texttt{DEBUG} builds}

If the \texttt{TEST} or \texttt{DEBUG} compile flags are set and
\texttt{zcash-ledger} is built, the resulting app is totally insecure. This is
by design to allow for testing: all on-device confirmations can be bypassed, and
secret keys can be replaced with fixed constants.

Some defenses should be added to the code to absolutely eliminate the
possibility of anyone accidentally side-loading a test build of
\texttt{zcash-ledger} onto their device. This could take the form of a warning
message that frequently gets displayed, e.g.\ prefix all of the normal text with
a ``WARNING: INSECURE TEST BUILD'' string.

\subsection{Compatibility testing with other wallets}

After Issue~\ref{issuezip32} is fixed, compatibility testing with other Zcash
shielded and transparent wallets should be implemented. This will prevent
accidental loss of funds in case any future bugs make \texttt{zcash-ledger}
incompatible with those other wallets.

\subsection{Increase compiler warning level} \label{recwarninglevel}

By default \texttt{zcash-ledger}'s build system does not pass any additional
\texttt{-W...} warning flags to the compiler. By adding \texttt{-Wall} or a
smaller custom set of warning flags, the compiler can be told to check for
various kinds of (likely) bugs. Once these warnings are fixed, we recommend
adding \texttt{-Werror} as well.

\subsection{Enforce stage-specific operations more strictly} \label{recmultiplefeeconfs}

To construct a transaction using \texttt{zcash-ledger}, the process proceeds
through a series of stages. The host first adds transparent inputs and outputs,
then Sapling outputs, then Orchard actions, then confirms the fee, and so on. At
each change of stage, the host must call the \texttt{CHANGE\_STAGE} APDU to
advance to the next stage.

Most of the crucial APDU commands, such as for adding transaction outputs, are
strictly enforced to occur during they stage in which they are expected.
However, several APDU commands, which could be stage specific, can be executed
by the host at any time. The state machine in \texttt{zcash-ledger} should be
hardened to disallow this behavior, i.e.\ all but the most general APDU commands
should be gated to a specific stage.

Doing this will place more requirement on the host to make calls in a specified
order, but will simplify the security analysis of the code.

Areas for hardening that we identified are:

\begin{enumerate}
      \item While in stage \texttt{FEE}, the host can call the
      \texttt{CONFIRM\_FEE} APDU to have the device calculate the SIGHASH and
      show the fee confirmation. The stage will remain in \texttt{FEE} until the
      user denies (stage returns to \texttt{IDLE}) or approves (stage advances
      to \texttt{SIGN}). While the user is confirming the fee, the stage is
      still in \texttt{FEE}, so the host can modify parts of the transaction
      (e.g.\ calling \texttt{SET\_S\_NET}) then call \texttt{CONFIRM\_FEE} again
      to show another fee confirmation. This behavior is unexpected, and should
      be ruled out.
      \item APDUs like \texttt{SET\_S\_NET}, \texttt{SET\_O\_NET},
      \texttt{SET\_*\_MERKLE\_PROOF}, and more can be called in any stage. While
      allowing extra flexibility makes developing the host wallet easier,
      strictly enforcing a stage for these APDUs would make it easier to analyze
      \texttt{zcash-ledger}'s code.
\end{enumerate}

In the case of (1), we were concerned that a race condition could potentially be
used by a malicious host to create an exorbitantly-large fee. The idea was to
(1) lie about the Sapling net value and call \texttt{CONFIRM\_FEE} so that the
fee confirmation looks normal but the transaction would be invalid, (2) while
the user is confirming the fee, call \texttt{SET\_S\_NET} again to provide the
correct Sapling net value, (3) call \texttt{CONFIRM\_FEE} again to recalculate
the SIGHASH. The hope was that a race condition could set the stage to
\texttt{SIGN} from the first fee confirmation \emph{after} the SIGHASH had been
recalculated with the exorbitant fee. This turns out not to work, but the reason
why depends on internal and possibly-changing details of how Ledger apps handle
APDUs. We recommend changing the code to make it very obvious this kind of
attack could never work. One way to do this is to have \texttt{CONFIRM\_FEE} set
the stage to a new \texttt{FEE\_CONFIRMING} stage immediately, so that
\texttt{CONFIRM\_FEE} cannot be called twice.

% ==================================================================
\section{Good Things} \label{secgoodthings}
% ==================================================================

\subsection{Fast security response}

In the process of \texttt{zcash-ledger}'s development, we reported several
issues to the author, all of which were fixed rapidly, often the same day. Some
of the issues were even noticed and fixed before we reported them.

\subsection{Integration testing}

Using the Ledger device simulator, \texttt{zcash-ledger} has integration tests
for most of its features. This is great for ensuring the app works as expected
and some of the tests would even catch certain kinds of cryptographic
implementation bugs.

% ==================================================================
\section{Future work} \label{secfuturework}
% ==================================================================

\subsection{Closer review of elliptic curve implementations} \label{sececcreview}

In this review, we did not thoroughly check the elliptic curve implementations.
The correctness of these implementations is crucial for security, since they are
used in the transaction signing process. Future security audits should review
them closely.

\subsection{More analysis of \texttt{zcash-ledger}'s state machine}

In the process of interacting with the \texttt{zcash-ledger} app, the host moves
the app through a series of stages. The \texttt{zcash-ledger} app informally
implements a state machine that moves through these stages.

State machine bugs are a common source of error in cryptographic software, and
\texttt{zcash-ledger}'s state machine's correctness is crucial to the app's
transaction confirmation process. We found two \emph{almost-bugs} in the state
machine (the issue in Appendix~\ref{issuestageadvance} and the almost-bug
described in Recommendation~\ref{recmultiplefeeconfs}). Therefore, we recommend
further review of \texttt{zcash-ledger}'s state machine.

% ==================================================================
\section{Conclusion} \label{secconclusion}
% ==================================================================

Let's summarize the issues again:

We found one \MediumIssue-severity issue that could put funds at risk:
side-channel defenses in the signing algorithm may be inadequate and could leak
information about spending keys to a malicious host.

We found several other \MediumIssue-severity issues: (a) because the Ledger
device's operating system does not support ZIP 32, the wallet is not
cross-compatible with other Zcash shielded wallets, (b) a malicious host can
``burn'' some of the user's funds by producing invalid ciphertexts, and (c) a
secure-UX issue where unified addresses are not displayed properly when the user
is confirming a transaction on the hardware device.

We also reported several \LowIssue-severity issues, which are mainly minor code
improvement suggestions, and pose little risk to users.

Many of the issues have simple fixes, but the incompatibility with ZIP 32 should
be a concern to the Zcash community. We either rely on Ledger to implement
proper ZIP 32 support in their operating system, or we need to complicate ZIP 32
itself with a workaround to support Ledger hardware wallets safely. If neither
of these paths succeed, we will be left with a confusing mess where wallets
derive keys from seeds differently, which will create a terrible user experience.

This audit \emph{did not} adequately review the elliptic curve implementations
in \texttt{zcash-ledger}. This has been left for future work.


\section*{Acknowledgements}

This security analysis was performed as part of the Zcash Ecosystem Security
Grant~\cite{zecsec, zecsecweb} funded by Zcash Community Grants~\cite{zcg}. Thanks to the
Zcash Grants Committee and the broader Zcash community for supporting the
security of privacy-enhancing open-source software.

Thanks to Daira Hopwood and Jack Grigg for guidance on how to remediate the ZIP
32 incompatibility issue.

\nocite{*}

\bibliographystyle{plainnat}
\bibliography{audit-report}

\pagebreak

\appendix
 
% ==================================================================
\renewcommand{\thesubsection}{\thesection.\arabic{subsection}}
\section{Bugs Discovered During Early Development} \label{appendixdevbugs}
% ==================================================================

In this section, we report some bugs that were found in pre-release versions of
the code. All of these issues were resolved prior to the first release (v1.0.1).

\subsection{Predictable nonces used while testing ECDSA signatures} \label{issueconstantnonce}
\textbf{Severity:} \CriticalIssue

In an \href{https://github.com/hhanh00/zcash-ledger/blob/7bcb456e32691acb2a2ea143c89e59a3d0b6f814/src/crypto/tx.c#L86}{in-development version of the code}, a constant of 32-bytes of repeated 3s was
used as the seed for the random number generator to generate the ECDSA nonces
for the shielded pool signatures.

This would have let the host predict the nonces, and so a malicious host could
have recovered the secret spending keys from the transaction's signatures.

This was fixed by generating the nonces at random using the device's secure
random number generator (except for test builds, which use predictable nonces to
enable testing).

\subsection{Fee confirmations could be bypassed by the host} \label{issueconfbypass}
\textbf{Severity:} \HighIssue

In an
\href{https://github.com/hhanh00/zcash-ledger/blob/6966e44d6ce9e2366eaeafcc6ccb18503f646347/src/crypto/tx.c#L375-L378}{in-development
version of the code}, the testing feature which allows bypassing confirmations
was not disabled in production builds for the fee confirmation.

This would have let a malicious host send arbitrary amounts of the user's funds
into the transaction fee. In cooperation with a miner, this would have allowed
for theft of funds.

It has been fixed by only allowing the confirmation bypassing to work in test
builds.

\subsection{Steal-by-fee by improperly advancing the \texttt{stage} state} \label{issuestageadvance}
\textbf{Severity:} \HighIssue

The host advances \texttt{zcash-ledger}'s state machine by calling the
\texttt{CHANGE\_STAGE} APDU, which calls the \texttt{change\_stage()} function:

\begin{lstlisting}[language=c]
int change_stage(uint8_t new_stage) {
    if (new_stage != G_context.signing_ctx.stage + 1) {
        reset_app();
        return io_send_sw(SW_BAD_STATE);
    }

    cx_hash_t *ph = (cx_hash_t *)&G_context.hasher;

    switch (new_stage) {
        case T_OUT:
        // ...
        // ... cases that handle stages S_OUT, O_ACTION, and FEE
        // originally there was no default case
    }
    G_context.signing_ctx.stage = new_stage;

    return io_send_sw(SW_OK);
}
\end{lstlisting}

The \texttt{new\_stage} parameter comes directly from the host. 

Originally, the switch statement did not have a \texttt{default} case to detect
bad state transitions. Therefore, a malicious host could:

\begin{enumerate}
      \item Wait for the user to initiate a legitimate transaction.
      \item Add extra inputs to the transaction so that it would have an enormous fee.
      \item Wait for the user to confirm all of the transaction outputs on the
      device. Now \texttt{stage} will be \texttt{FEE}.
      \item Initiate the fee confirmation, which computes the transaction's
      SIGHASH for signing and shows the fee to the user, but \texttt{stage}
      should not change to \texttt{SIGN} unless the user approves.
      \item Call \texttt{CHANGE\_STAGE} with \texttt{new\_stage = SIGN = FEE +
      1} before the user rejects the bad fee. Now \texttt{stage} is
      \texttt{SIGN} and signatures for the transaction can be extracted by
      calling the signing APDUs.
\end{enumerate}

In other words, a malicious host could have sent arbitrary amounts of the user's
funds into a transaction fee if the user was not quick enough to disapprove of
the displayed bad fee. In cooperation with a miner, this could have been used to
steal funds.

This was fixed by adding a \texttt{default} case that rejects bad state
transitions:

\begin{lstlisting}[language=c]
      // ...
        default: // should have been caught by the check at the top
            return io_send_sw(SW_INVALID_PARAM);
      // ...
\end{lstlisting}

\subsection{Out-of-bounds memory access in the FF1 implementation} \label{issueoob}
\textbf{Severity:} \LowIssue

In an old version of the FF1 implementation, a \texttt{memmove} made an
out-of-bounds write:

\begin{lstlisting}[language=c]
int ff1(uint8_t *d, const uint8_t *dk, uint8_t *di) {
      // ...
      memset(P, 0, 16);
      P[9] = i;
      memmove(&P[10], b, 6);
      // block depends on the round number and b
      error = cx_aes_iv_no_throw(&aes_key, CX_ENCRYPT | CX_PAD_NONE | CX_CHAIN_CBC | CX_LAST, (uint8_t *)R, 16, 
          (uint8_t *)P, 16, (uint8_t *)R, &out_len);
      if (error) return error;
      memmove(d, R, 16);
      // ...
\end{lstlisting}

The \texttt{d} written to in the last line was allocated with a size of only 11
bytes in \texttt{orchard.c}, so the 16-byte write would overflow:

\begin{lstlisting}[language=c]
// ...
uint8_t d[11];
memset(d, 0, 11);
ff1(d, G_context.orchard_key_info.dk, d);
// ...
\end{lstlisting}

This code can be found in commit \texttt{afc3a2b0b0007639a865b9415470591abbecd17b}.
It has been fixed by entirely removing the \texttt{ff1()} function, replaced by
\texttt{ff1\_inplace()} which does not appear to suffer from the same issue.

\end{document}